\chapter*{Note on Epistemic Status}
\addcontentsline{toc}{chapter}{Note on Epistemic Status}

This monograph distinguishes, throughout and without exception, between
results that are proved, results that are conjectured, questions that
are left explicitly open, and claims that are stated informally, ahead
of formalization, as a working thesis rather than as any of the above.
Five labels are used accordingly. A \emph{Theorem} is a claim with a
complete proof given in the chapter or in Appendix~\ref{app:proof}. A
\emph{Proposition} is a claim with a complete, if sometimes short,
proof, of lesser generality or centrality than a Theorem. A
\emph{Conjecture} is a claim believed true, motivated by argument or by
worked examples, but not proved in this monograph. An \emph{Open
Question} is a question the monograph raises and explicitly declines to
resolve, together with a statement of what resolving it would require.
A \emph{Principle} is a pre-formal statement of a thesis --- most
importantly the informal statement of global insufficiency in
Chapter~\ref{ch:continuation-problem} and its restatement in
Chapter~\ref{ch:synchronic-to-diachronic} --- introduced before the
vocabulary needed to state it precisely has been built, and superseded
once a formal Proposition or Theorem makes the same claim rigorously;
it carries no claim to proof of its own.

Two results carry the label \emph{Theorem}: the Redistribution Theorem
(\cref{thm:redistribution}), imported without modification from
\emph{Conservation of Ambiguity} and explicitly scoped there as a
synchronic, single-observer result; and the Composition Failure
Theorem (\cref{thm:composition-failure}), the new diachronic,
multi-agent result this monograph exists to establish. Every other
theorem-like claim in the book is labeled Proposition, Corollary,
Conjecture, or Open Question according to the criteria above, and
Appendix~\ref{app:open-questions} consolidates every item labeled
Conjecture or Open Question in one place, together with every
substantive open item flagged in a Remark rather than a formal
environment, so that the book's unproved commitments are never left
implicit.

The single largest such commitment is the status of \emph{MEM|8}'s
relationship between memory and representation (\cref{ch:mem8}), which
this monograph treats as a genuine, unresolved fork rather than
resolving by assumption. Readers should treat Part~\ref{part:modes-of-access}'s
results as conditional on the narrower of the two available readings
until that fork is settled by further work.
